%------------------------------------------------------------------------------%
% LINEAR ALGEBRA                                                               %
%------------------------------------------------------------------------------%
% File  : Linear_Algebra.tex                                                   %
% Author: Klaus Hoeflinger, Beethovenstrasse 11, D-73117 Wangen                %
% Email : khoeflinger@t-online.de                                              %
%------------------------------------------------------------------------------%

%------------------------------------------------------------------------------%
% DOCUMENT CLASS and INCLUDE FILES                                             %
%------------------------------------------------------------------------------%
\documentclass[11pt,twoside]{article}
\usepackage{geometry}
\geometry{paperheight=241mm,
          paperwidth=152mm,
          top=22mm,
          bottom=17mm,
          left=20mm,
          right=20mm,
          textwidth=112mm,
          textheight=202mm,
          headheight=10mm,
          headsep=3mm,
          footskip=9mm,
          includehead,
          includefoot,
          marginparsep=0mm,
          marginparwidth=0mm}

\renewcommand{\baselinestretch}{0.945}\normalsize

\AtBeginDocument{\setlength\abovedisplayskip{2mm}}
\AtBeginDocument{\setlength\belowdisplayskip{2mm}}

%------------------------------------------------------------------------------%
% define title formats                                                         %
%------------------------------------------------------------------------------%
\usepackage{titlesec}

\renewcommand{\thesection}{\arabic{section}}
\renewcommand{\sfdefault}{FiraSans}
\renewcommand{\ttdefault}{pcr}
\renewcommand{\rmdefault}{antiqua}

\titleformat{\part}[display]
{\normalfont \large \rmfamily \bfseries \centering}
{\small{\thepart.}}{2pt}{}

\titleformat{\section}[runin]
{\normalfont \rmfamily \bfseries}
{\thesection}{1pt}{.\;}[.]

%\titleformat{\section}
%{\normalfont \bfseries \small \filcenter}
%{\thesection.\;}{0mm}{}

\titleformat{\subsection}[runin]
{\normalfont \itshape}
{}{0mm}{}

\titleformat{\subsubsection}[runin]
{\normalfont \itshape}
{}{}{}{}

\titleformat{\paragraph}[runin]
{\normalfont}
{}{}{}{}

\titleformat{\subparagraph}[runin]
{\normalfont}
{}{}{}{}

\titlespacing*{\part}          { 0pt}{ 0pt}{ 8pt}
\titlespacing*{\section}       { 0pt}{ 7pt}{ 4pt}
\titlespacing*{\subsection}    { 0pt}{14pt}{ 6pt}
\titlespacing*{\subsubsection} { 0pt}{ 6pt}{12pt}
\titlespacing*{\paragraph}     { 0pt}{ 6pt}{12pt}
\titlespacing*{\subparagraph}  { 0pt}{ 0pt}{12pt}

%------------------------------------------------------------------------------%
% define enumerate parameters                                                  %
%------------------------------------------------------------------------------%
\usepackage{enumitem}
\setlength{\parindent}{3pt}
\setlength{\parskip}  {0pt}

%\setlist{noitemsep}

\setlist[enumerate]{
         leftmargin=12pt,
         rightmargin=0pt,
         itemsep=1pt,
         itemindent=3pt,
         topsep=3pt,
         labelsep=4pt,
         partopsep=0pt,
         labelwidth=0pt,
         listparindent=0pt,
         parsep=0pt}

\setlist[itemize]{
         leftmargin=12pt,
         rightmargin=0pt,
         itemsep=0pt,
         itemindent=1pt,
         topsep=1pt,
         labelsep=4pt,
         partopsep=1pt,
         labelwidth=0pt,
         listparindent=0pt,
         parsep=0pt}

%------------------------------------------------------------------------------%
% define packages                                                              %
%------------------------------------------------------------------------------%
\usepackage{upref}
\usepackage{amssymb}
\usepackage{slantsc}
\usepackage{amsmath}
\usepackage{amsthm}
\usepackage{thmtools}
\usepackage{amscd}
\usepackage{verbatim}
\usepackage{latexsym}
\usepackage{multicol}
\usepackage{graphicx}
\usepackage{setspace}
\usepackage[ngerman,english]{babel}
\usepackage[protrusion=true,expansion=true]{microtype}
\usepackage{tikz}
\usetikzlibrary{arrows,arrows.meta,decorations.markings}
\usepackage{mathrsfs}
\usepackage{amsfonts}
\usepackage{datetime2}
\usepackage{textgreek}
\usepackage{hyperref}
\usepackage{pifont}
\usepackage{relsize}
%\usepackage{biblatex}
\relsize{-0.05}
\usepackage[skip=0mm]{parskip}
\usetikzlibrary{decorations.pathmorphing}

\usepackage{mlmodern}
\usepackage[T5,T1]{fontenc}
\usepackage{ragged2e}

%\usepackage{mathabx}
%\usepackage{times}

%\usepackage{fontspec}
%\setmainfont{Nimbus Roman No9 L}
%\usepackage[T1]{fontenc}

%------------------------------------------------------------------------------%
% define page layout                                                           %
%------------------------------------------------------------------------------%
\usepackage{fancyhdr}
\pagestyle{fancy}
\setlength{\headheight}{30.0pt}
\renewcommand{\headrulewidth}{0pt}

\AtBeginDocument{
  \addtolength\abovedisplayskip{-0.0\baselineskip}
  \addtolength\belowdisplayskip{-0.0\baselineskip}
}

\fancyhead{}
\fancyfoot{}
\addtolength{\headwidth}{0.02\marginparwidth}

\makeatletter
\let\ps@plain\ps@fancy
\makeatother

\renewcommand\sectionmark[1]
{
 \def\sectionname{#1}
 \markright{\thesection #1}
}

\makeatletter
\@addtoreset{section}{part}
\makeatother

%------------------------------------------------------------------------------%
% define footnote without number                                               %
%------------------------------------------------------------------------------%
\newcommand{\myfootnote}[1]{%
\renewcommand{\thefootnote}{}%
\footnotetext{#1}%
\renewcommand{\thefootnote}{\arabic{footnote}}%
}

%------------------------------------------------------------------------------%
% define numbering in equations                                                %
%------------------------------------------------------------------------------%
\numberwithin{equation}{section}

%------------------------------------------------------------------------------%
% define newtheorem                                                            %
%------------------------------------------------------------------------------%
\declaretheoremstyle[
headfont=\scshape, 
headindent = 0mm, %\parindent,
%postheadhook = {\hspace{0mm}\newline},
%postheadspace = \newline, 
spaceabove = 3mm, 
spacebelow = 3mm,
numberwithin=section,
name=Theorem]{khMATH}
\declaretheorem[style=khMATH]{theorem}

\declaretheoremstyle[
headfont=\bfseries, 
headindent = 0mm, %\parindent,
%postheadhook = {\hspace{0mm}\newline},
%postheadspace = \newline, 
spaceabove = 3mm, 
spacebelow = 3mm,
numberwithin=part,
name=Theorem]{khMATH1}
\declaretheorem[style=khMATH1]{theorem1}

\declaretheoremstyle[
headfont=\scshape, 
headindent = 0mm, %\parindent,
%postheadhook = {\hspace{0mm}\newline},
%postheadspace = \newline, 
spaceabove = 3mm, 
spacebelow = 3mm,
numberwithin=section,
name=Corollary]{khMATH1}
\declaretheorem[style=khMATH1]{corollary}

%            name         counter     label              numeration-mode
\theoremstyle{plain}
\newtheorem*{introduction}            {\sc Introduction}
\newtheorem {standard}                {}                 [section]
\newtheorem*{definition}              {\sc Definition}
\newtheorem{satz1}                    {\sc Satz}
\newtheorem{lemma}                    {\sc Lemma}
\newtheorem*{proposition}             {\sc Proposition}
%\newtheorem{corollary}               {\sc Corollary}
\newtheorem*{corollar}                {\sc Korollar}
\newtheorem*{remark}                  {\sc Remark}
\newtheorem*{remarks}                 {\sc Remarks}
\newtheorem*{example}                 {Example}
\newtheorem*{examples}                {Examples}
\newtheorem*{theo}                    {\sc Theorem}

%------------------------------------------------------------------------------%
% define tabular                                                               %
%------------------------------------------------------------------------------%
\usepackage{tabularx}
\newcolumntype{L}[1]{>{\raggedright\arraybackslash}p{#1}}
\newcolumntype{C}[1]{>{\centering\arraybackslash}  p{#1}} 
\newcolumntype{R}[1]{>{\raggedleft\arraybackslash} p{#1}} 

%------------------------------------------------------------------------------%
% define \sh, \zB                                                              %
%------------------------------------------------------------------------------%
\def\dsh{{d.\,h.}}
\def\ie{{i.\,e.}}
\def\zB{z.\,B.}
\renewcommand{\abstractname}{ }

%------------------------------------------------------------------------------%
% change default spacing for cases                                             %
%------------------------------------------------------------------------------%
\makeatletter
\renewcommand*\env@cases[1][1.6]{%
  \let\@ifnextchar\new@ifnextchar
  \left\lbrace
  \def\arraystretch{#1}%
  \array{@{}l@{\quad}l@{}}%
}
\makeatother

\newenvironment{rcases}
  {\left.\begin{aligned}}
  {\end{aligned}\right\rbrace}

%------------------------------------------------------------------------------%
% change default for \predisplaypenalty                                        %
%------------------------------------------------------------------------------%
\predisplaypenalty=990
\allowdisplaybreaks \numberwithin{equation}{section}

%------------------------------------------------------------------------------%
% general notations                                                            %
%------------------------------------------------------------------------------%
\def\*{\star}
\def\={\,=\,}
\def\+{\,+\,}
\def\xsubset{{\,\subset\,}}
\def\xtimes{{\,\times\,}}
\def\xeof{{\,\in\,}}
\def\eq{{\,=\,}}
\def\xto{{\,\to\,}}
\def\eqdef{\,:=\,}
\def\:{{\,:\,}}
\def\ele{{\,\in\,}}
\def\exists{\mbox{ exists }}
\def\and{\mbox{ and }}
\def\for{\mbox{ for }}
\def\fa{\mbox{ for all }}
\def\fe{\mbox{ for each }}
\def\qfa{\quad \mbox{ for all }}
\def\df{\,\dotfill}
\def\iff{if and only if }
\def\wrt{with respect to}

\makeatletter
\newcommand \Df {\leavevmode \cleaders \hb@xt@ .70em{\hss .\hss }\hfill \kern \z@}
\makeatother

\def\sql{\mathfrak{a}}               % sesquilinear form a
\def\DF{B}                           % Dirichlet form a
%------------------------------------------------------------------------------%
% number sets and special elements                                             %
%------------------------------------------------------------------------------%
\def\P{\mathbb{H}}                   % half plane of $\C$
\def\J{I}                            % interval I
\def\N{{\mathbb{N}}}                 % unsigned integers
\def\Q{{\mathbb{Q}}}                 % rational integers
\def\Z{{\mathbb{Z}}}                 % integers
\def\R{{\mathbb{R}}}                 % real numbers
\def\Rn{{\mathbb{R}}^n}              % real numbers in Rn
\def\Rd{{\mathbb{R}}^d}              % real numbers in Rn
\def\C{{\mathbb{C}}}                 % complex numbers
\def\F{{\mathbb{K}}}                 % arbitrary field
\def\Torus{{\mathbb{T}}}             % complex circle line
\def\T{\tau}                         % upper bound of interval (0,t)
\def\sign#1{\operatorname{sign}(#1)} % sign of a number

%------------------------------------------------------------------------------%
% set theory                                                                   %
%------------------------------------------------------------------------------%
\def\G{G}                            % domain $G$
\def\clG{\overline{\G}}              % closure of $G$
\def\bndG{\partial \G}               % boundary of $G$
\def\setA{\mathcal{A}}               % set A
\def\setB{\mathcal{B}}               % set B
\def\setC{\mathcal{C}}               % set C
\def\setM{M}                         % set M
\def\setN{N}                         % set N
\def\setS{\mathcal{S}}               % set S
\def\famF{\mathfrak{F}}              % family of sets
\def\indI{\mathcal{I}}               % index set I
\def\pot#1{\mathfrak{P}(#1)}         % power set

\def\relR{\mathbf{R}}                % relation R
\def\mapf{f}                         % mapping f
\def\mapg{g}                         % mapping g

\def\set#1#2{\{ \, #1 \,: \, #2 \, \}}
\def\tensor#1#2{#1 \otimes #2}
\def\Tens{\oplus}

\def\cal#1{{\mathcal{#1}}}
\def\aut#1{\operatorname{Aut}(#1)}

\def\cross#1#2{#1 {\,\times\,} #2}
\def\multicross#1#2{#1 \times  \ldots \times  #2}
\def\multitens#1#2 {#1 \otimes \ldots \otimes #2}

%------------------------------------------------------------------------------%
% group theory                                                                 %
%------------------------------------------------------------------------------%
\def\1{{\mathsf{1}}}                % unit element of a monoid
\def\0{{\mathsf{0}}}                % unit element of an Abelian monoid
\def\kernel#1{\mathfrak{N}(#1)}     % kernel of a homomorphism
\def\cokernel#1{\mathfrak{C}{(#1)}} % cokernel of a homomorphism
\def\Hom#1{\operatorname{Hom}(#1)}

%------------------------------------------------------------------------------%
% linear algebra                                                               %
%------------------------------------------------------------------------------%
\def\vecV{\mathit{V}}               % vector space V
\def\vecH{\mathit{H}}               % vector space H
\def\vecW{\mathit{W}}               % vector space W
\def\vecU{\mathit{U}}               % subspace U
\def\convC{\mathcal{C}}             % convex set C
\def\basH{\mathfrak{B}}             % Hamel base B

\def\LO#1{\mathscr{L}(#1) }         % algebra of linear transformations

\def\scalar#1#2{ \langle #1,#2 \rangle }
\def\trace#1{\operatorname{tr}(#1)}
\def\dim#1{{\operatorname{dim} \, #1 }}
\def\codim#1{{\operatorname{codim} \, #1 }}
\def\dim{\operatorname{dim}}
\def\codim{\operatorname{codim}}
\def\ind{\operatorname{ind}}

\def\genG{\cal{G}}

\def\algA{\cal{A}}
\def\algB{\cal{B}}
\def\algU{\cal{U}}

\def\idI{\cal{I}}

%------------------------------------------------------------------------------%
% topology                                                                     %
%------------------------------------------------------------------------------%
\def\compK{{K}}          % compact set C
\def\openO{\mathcal{O}}  % open set O
\def\openU{\mathcal{U}}  % open set U
\def\U{\mathcal{U}}      % neighbourhood U
\def\V{\mathcal{V}}      % neighbourhood V
\def\d{\mathit{d}}       % metric function d

\def\interior#1{\mathring{#1}}
\def\closure#1{{\overline{#1}}}

\def\topT{\mathcal{T}}   % topology T
\def\topX{X}             % topological space X
\def\topY{Y}             % topological space Y
\def\topZ{Z}             % topological space Z
\def\metrS{M}            % metric space M

\def\grpG{G}             % topological group G
\def\grpA{A}             % topological group A
\def\grpB{B}             % topological group B
\def\grpC{C}             % topological group C
\def\grpH{H}             % topological group H
\def\grpU{U}             % topological group U
\def\grpV{V}             % topological group V
\def\topG{G}             % topological group G
\def\topH{H}             % topological group H
\def\topU{U}             % topological subgroup U

\def\subS{\cal{S}}
\def\riR{R}              % ring R

%------------------------------------------------------------------------------%
% calculus                                                                     %
%------------------------------------------------------------------------------%
\def\supp#1{\operatorname{supp}(#1)}
\def\singsupp#1{\operatorname{sing \, supp}(#1)}

\def\re{\operatorname{Re}}
\def\im{\operatorname{Im}}
\def\grad{\operatorname{grad}}

%\def\abs#1 {\left|#1\right|}
\def\abs#1 {\lvert #1 \rvert}
%\def\abs#1 {\left|\,#1\,\right|}

%------------------------------------------------------------------------------%
% measure theory                                                               %
%------------------------------------------------------------------------------%
\def\salgA{\Sigma}               % sigma-algebra S
\def\salgB{\cal{B}}              % sigma-algebra B
\def\Borel#1{\cal{B}_{#1}}       % Borel-algebra 
\def\LB{\lambda}                 % Lebesgue-Borel measure

%------------------------------------------------------------------------------%
% function spaces                                                              %
%------------------------------------------------------------------------------%
\def\Lp{L^p(\G)}                 % spaces L_p
\def\Lq{L^q(\G)}                 % spaces L_q
\def\Lh{L^2(\G)}                 % spaces L_2
\def\Li{L^{\infty}(\G)}          % spaces L_infty
\def\Lc{L_{loc}^1(\G)}           % spaces L_1^{loc}
\def\Ck{C^k(\G)}                 % spaces C_k

\def\BV#1{BV(#1) }               % set of functions of bounded variation
\def\AC#1{AC(#1) }               % set of absolutely continuous functions

\def\MP#1{\mathscr{M}^+(#1) }    % set of positive measures
\def\MS#1{\mathscr{M}^-(#1) }    % vector space of finite signed measures
\def\MC#1{\mathscr{M}(#1) }      % vector space of complex measures
\def\MCR#1{\mathscr{M}_{r}(#1) } % vector space of regular complex measures

\def\SobW{W^{k,p}(\G)}                % Sobolev space W
\def\SobO{\mathring{W}^{k,p}(\G)}     % Sobolev space W
%\def\WN#1#2{\mathring{W}^{#1}(#2) }   % Sobolev space W
\def\WN#1#2{W_0^{#1}(#2) }   % Sobolev space W
%\def\HN#1#2{\mathring{H}^{#1}(#2) }   % Sobolev space H
\def\HN#1#2{H_0^{#1}(#2) }   % Sobolev space H
%------------------------------------------------------------------------------%
% functional analysis                                                          %
%------------------------------------------------------------------------------%
\def\snorm#1{\lVert #1 \rVert}
\newcommand{\norm}[1]{\left\lVert #1 \right\rVert}

\def\topV{V}                       % topological vector space V
\def\topW{W}                       % topological vector space W
\def\normS{X}                      % normed space

\def\banX{\mathit{X}}              % Banach space X
\def\banY{\mathit{Y}}              % Banach space Y
\def\banZ{\mathit{Z}}              % Banach space Z
\def\dbanX{\mathit{X^*}}           % dual space of Banach space X
\def\dbanY{\mathit{Y^*}}           % dual space of Banach space Y
\def\dbanZ{\mathit{Z^*}}           % dual space of Banach space Z

\def\banA{\mathit{A}}              % Banach algebra A
\def\banB{\mathit{B}}              % Banach algebra B
\def\banC{\mathit{C}}              % C*-algebra C


\def\domain#1{\mathfrak{D}(#1)}    % domain
\def\range#1{\mathfrak{R}(#1)}     % range

\def\isoJ{\mathfrak{J}}            % isometry J
\def\repA{{\cal{A}_{\sql}}}        % representation operator A
\def\linS{{S}}                     % linear operator S
\def\linG{{G}}                     % linear operator G
\def\linL{{L}}                     % linear operator L
\def\linT{{T}}                     % linear operator A
\def\linA{\mathit{A}}              % linear operator A
\def\linB{{B}}                     % linear operator B
\def\linM{{M}}                     % linear operator M
\def\linU{{U}}                     % linear operator U
\def\linI{{I}}                     % linear operator I
\def\A{\cal{A}}                    % linear differential operator A
\def\BDO{\cal{B}}                  % linear boundary differential operator B
\def\linU{{U}}                     % unitary operator U
\def\linP{{P}}                     % projection P
\def\compA{k}                      % compact operator k
\def\linFR{f}                      % finite rank operator f
\def\fredA{A}                      % Fredholm operator F

\def\HAM{\cal{H}}                  % Hamilton operator
\def\PA{\partial^{\alpha}}         % elementary differential operator
\def\DO{\cal{A}}                   % linear differential operator
\def\BC{\cal{B}}                   % boundary operator
\def\SLO{\cal{L}_{p,q}}            % Sturm-Liouville operator
\def\MO{\cal{L}_{M}}               % momentum operator

\def\HS#1{S(#1)           }        % algebra of Hilbert-Schmidt operators
\def\FR#1{F(#1)           }        % algebra of finite rank operators
\def\TC#1{N(#1)           }        % algebra of trace class operators
\def\BU#1{\mathscr{U}(#1) }        % algebra of unitary linear operators
\def\BO#1{\mathscr{L}(#1) }        % algebra of bounded linear operators
\def\CO#1{\mathscr{C}(#1) }        % algebra of compact linear operators
\def\CL#1{\mathscr{C}(#1) }        % algebra of compact linear operators
\def\FO#1{\Phi(#1) }               % set of Fredholm operators
\def\FK#1#2{\Phi_{#2}(#1) }        % set of Fredholm operators with index k

\def\hilH{\mathit{H}}              % Hilbert space H
\def\clU{\mathit{U}}               % closed subspace U

\def\orthO{\mathfrak{O}}           % orthonormal system O
\def\orthS{\mathfrak{S}}           % orthonormal system S
\def\basB{\mathfrak{B}}            % basis B of a Hilbert space

\def\GL#1 {\mathbf{GL}(#1)}        % linear group
\def\OG#1 {\mathbf{O}(#1)}         % linear group
\def\SOG#1 {\mathbf{SO}(#1)}       % linear group

%------------------------------------------------------------------------------%
% distribution theory                                                          %
%------------------------------------------------------------------------------%
\def\disF{F}                  % distribution F
\def\disG{G}                  % distribution G
\def\disH{H}                  % distribution H
\def\disU{U}                  % distribution U
\def\disT{T}                  % distribution T
\def\SF{C^{\infty}(\G)}       % space of smooth functions
\def\TF{C_c^{\infty}(\G)}     % space of compactly supported smooth functions
\def\TFRd{C_c^{\infty}(\R^d)} % space of test functions
\def\SRd{\mathscr{S}(\R^d)}   % Schwartz space
\def\SR1{\mathscr{S}(\R)}     % Schwartz space
\def\B1#1{C^1(#1) }           % space of continuously differentiable functions
\def\Bm#1{C^m(#1) }           % space of continuously differentiable functions
\def\Ck{C^k(\G)}
\def\TD{\mathscr{S}'(\R^d)}   % space of temperated distributions
\def\E{\cal{E}}               % space of distributions with compact support
\def\D{\mathscr{D}}           % D
\def\FT{\mathcal{F}}          % Fourier transformation F
\def\FT{\mathscr{F}}          % Fourier transformation F

%------------------------------------------------------------------------------%
% other definitions                                                            %
%------------------------------------------------------------------------------%
\def\Proof{\textsc{Proof}}
\def\FRECHET{Fr\'{e}chet}
\def\ARZELA{Arzel\'{a}}
\def\GARDING{G\r{a}rding}
\def\HOELDER{H\"older}
\def\SCHROEDINGER{Schr\"odinger}
\def\JOERGENS{J\"orgens}
\def\WUEST{W\"uest}
\def\KAEHLER{K\"aehler}
\def\HOEFLINGER{H\"oflinger}
\def\POINCARE{Poincar\'{e}}
\def\FA{Functional Analysis}

\def\Author   {K.\,{\HOEFLINGER}}
\def\AuthorUC {K.\,HOEFLINGER}
\def\Version  {1.0}
\def\Email    {khoeflinger@t-online.de}


\titleformat{\part}%[display]
{\normalfont \rmfamily \bfseries \centering}
{{\thepart.}}{3pt}{}

\titleformat{\section}[runin]
{\normalfont \bfseries \filcenter}
{\thesection.\,}{0mm}{}[.]

\titlespacing*{\part}   {5pt}{12pt}{ 4pt}
\titlespacing*{\section}{0pt}{ 3pt}{ 3pt}

%\renewcommand\thesection{\Roman{section}}
%\renewcommand*{\thesubsection}{\arabic{subsection}}

%------------------------------------------------------------------------------%
%                                BEGIN OF DOCUMENT                             %
%------------------------------------------------------------------------------%
\begin{document}
\thispagestyle{empty}

\centerline{\textrm{\large{\textbf{Linear Algebra}}}}
\vspace{4pt}
\centerline{\textsc{\small{\Author}}}
\vspace{10pt}

\fancyhead[OL]{\small{\thepage}}
\fancyhead[ER]{\small{\thepage}}
\fancyhead[OC] {\small{\textsc{Linear Algebra}}}
\fancyhead[EC]{\small{\thesection}. \textsc{\sectionname}}


%------------------------------------------------------------------------------%
%                                  CONTENT                                     %
%------------------------------------------------------------------------------%
\myfootnote
{
  Version {\Version} from \today;
  Email:  \textit{\Email}
}

\textit{Linear algebra} is the mathematical subject of modules and linear spaces
together with the linear mappings defined between them.
It is the fundament for a couple of other topics in mathematics,
which in fact relies on the concepts provided by linear algebra.
The content is divided up into three main parts,
namely \textit{rings and fields}, \textit{linear spaces and mappings}
and finally \textit{algebras}.

\paragraph{}
Let us first recall from group theory
that an \textit{Abelian group} is a set $\grpG$ together with
an associative operation $+ \:\grpG \times \grpG \xto \grpG$
such that

\begin{itemize}[leftmargin=15pt,itemsep=2pt]
\item[1.]
there is exactly one element $\0 \xeof \grpG$
such that the computation rule
$x + \0 \eq \0 + x \eq x$ holds for all other $x \xeof \grpG$,
and $\0$ is called the \textit{unit element} in $\grpG$.
\item[2.]
to each $x \xeof \grpG$ there is a unique element $-x \xeof \grpG$,
called the \textit{inverse} of $x$
such that $x + (-x) \eq -x + x \eq \0$.
\item[3.]
the operation $+$ is \textit{commutative} or \textit{Abelian},
{\ie} $x+y \eq y+x$ for all $x,y \xeof \grpG$.
\end{itemize}

\setcounter{page}{1}
%------------------------------------------------------------------------------%
\fancyhead[EC] {\small{I. \textsc{Rings and Fields}}}
\fancyhead[OC] {\small{I. \textsc{Rings and Fields}}}

\part{Rings and Fields}  
%------------------------------------------------------------------------------%
The concepts of \textit{ring} and \textit{field}
constitute algebraic structures,
which form the base for the domains of numbers $\Z, \Q, \R$ and $\C$.
Moreover,
rings and fields are used to define further structures
like \textit{modules}, \textit{vector spaces} and \textit{algebras}.

\paragraph{}
In the following let $(R,+)$ be an Abelian group
and $\circ \:R \times R \xto R$ be a further associative operation.
The triple $(R,+,\circ)$ is called a \textit{ring},
if the operations $+$ and $\circ$ satisfy the following computation rules,
denoted as \textit{distributive laws}:
\[
\begin{cases}[1.4]
\quad x \circ (a + b) = (x \circ a) + (x \circ b), \\
\quad (a + b) \circ y = (a \circ y) + (b \circ y), \\
\end{cases}
\]
for $a,b,x,y \xeof R$.
The pair $(R,\circ)$ especially forms a semigroup,
and the operation $\circ$ is usually called the \textit{ring multiplication} in $R$.

\subparagraph{}
Any ring $R$ is called a \textit{ring with unit element},
if there is a element $1 \xeof R$ such that $x \circ \1 \eq \1 \circ x \eq x$
for all $x \xeof R$.
Moreover,
a ring is called \textit{commutative},
if the multiplication is so,
that is,
$(R,\circ)$ constitutes a commutativee semigroup.

\paragraph{}
The following examples of rings are well-known:

\begin{itemize}[leftmargin=12pt]
\item[1.]
$(\Z,+,\cdot)$ with unit element $1 \xeof \Z$ is a commutative ring.

\item[2.]
Let $\setM$ be a set, $(R,+,\circ)$ be a ring
and $\cal{A}(M,R)$ be the set of mappings from $\setM$ to $R$.
We introduce addition and multiplication of mappings
$f,g \xeof \cal{A}(M,R)$ according to
\[
\begin{cases}[1.4]
\quad (f + g)(x) := f(x) + g(x),\\
\quad (f \cdot g)(x) := f(x) \cdot g(x),\\
\end{cases}
\]
for all $x \xeof M$.
The set $(\cal{A}(M,\R), +, \cdot)$ forms a ring as well.

\item[3.]
Let $(A,+)$ be an Abelian group
and $\cal{A}(A)$ be the set of mappings $f \: A \xto A$.
Let for $f,g \xeof \cal{A}(A)$ the sum $f + g$ and the prodzct
$f \circ g$ be given by
\[
(f + g)(x) := f(x) + g(x), \quad (f \circ g)(x) := f(g(x))
\]
for all $x \xeof A$.
Then $(\cal{A}(A),+,\circ)$ is a non-commutative ring with unit element $\1$,
where $\1$ is nothing but the identity mapping in $A$.

\item[4.]
The set of $(n,n)$-matrices with coefficients in $\R$ or 
$\C$ forms a ring with unit element {\wrt} matrix multiplication,
which is non-commutative for $n > 1$.
\end{itemize}

\paragraph{}
Let us turn to some subclasses of rings
by posing some additional requirements regarding the ring multiplication:

\subparagraph{}
Let $(R,+,\circ)$ be a ring.
An element $r \xeof R$ is called a \textit{zero divisor} of $R$,
if $r \neq 0$ and there is $x \xeof R$
such that $x \neq 0$ and $r \circ x \eq 0$.
A ring $R$ is called \textit{without zero divisors},
if it does not contain a zero divisor.

\subparagraph{}
The absence of a zero divisor in a ring can be characterized
by means of the \textit{reduction rule}:
a ring is without zero divisors
\iff
for $x,y,z \xeof R$ and $x \neq 0$ the equation $x \circ y \eq x \circ z$
implies $y \eq z$.
Damit gelangen we zum Begriff des \textit{Integrit\"{a}ts ring}:
a ring $R$ is called a \textit{Integrit\"{a}tsring},
if it is without zero divisors and commutative 
and owns a unit element $1 \neq 0$.

\subparagraph{}
Let $(R,+,\circ)$ be a ring with unit element $1 \xeof R$.
An element $r \xeof R$ is called a \textit{unit} rsp. \textit{invertible},
if there is $r' \xeof R$ such that $r \circ r' \eq 1$.
If $r'$ exists, then it is uniquely defined by $r$
(notation: $r' \eq r^{-1}$).
If $r$ and $r^{-1}$ are both units, then $(r^{-1})^{-1} \eq r$ holds.

\subparagraph{}
We denote $E(R)$ to be the set of units of a ring $R$.
The pair $(E(R),\circ)$ forms an Abelian group for every ring $R$.

\section{Subrings}
%------------------------------------------------------------------------------%
Subrings are defined in a smiliar way as subgroups.
Let $(R,+,\circ)$ be a ring.
Then yny set $U \subset R$ is called a \textit{subring} of $R$, 
if for all $r,s \xeof U$ we have $r+s \xeof U$ and $r \circ s \xeof U$.

\subparagraph{}
Hence a subring is closed {\wrt} the algebraic operations $+$ and $\circ$.
Each subring $(U,+)$ is an (Abelian) subgroup of $(R,+)$,
and $(U,\circ)$ is a sub-semigroup of $(R,\circ)$.
The intersection of arbitrarily many subrings of a ring is a subring again.
If $R$ is a ring and $S \subset R$ is a non-empty subset,
then the \textit{generated subring} $[S]$ of $S$ is defined
to be the intersection of all subrings $U$ of $R$
having the property $S \xsubset U$.

\section{Ideals}
%------------------------------------------------------------------------------%
One of the most important concepts in ring theory is that of \textit{ideal}.
Let $(R,+,\circ)$ be a ring and $I \subset R$ be a subring:

\begin{itemize}[leftmargin=15pt]
\item[1.]
$I$ is called a \textit{left-sided ideal} in $R$,
if for each $r \xeof R$ and each $x \xeof I$ is $r \circ x \xeof I$.
\item[2.]
$I$ is called a \textit{right-sided ideal} in $R$,
if $x \circ r \xeof I$ for each $r \xeof R$ and each $x \xeof I$.
\item[3.]
$I$ is called \textit{two-sided},
if it is a left and right-sided ideal at the same time.
\end{itemize}

\subparagraph{}
It follows from this definition
that for ideals in a \textit{commutative} ring all this notations
are equivalent to each other.

\subparagraph{}
The intersection of arbitrarily many ideals in $R$ forms an ideal again.
Hence we can assign to each non-empty subset $S \subset R$ a \textit{smallest} ideal,
wich contains the set $S$.
The ideal generated by the set $S$ is denoted $[S]$ and is defined
as the intersection of all ideals in $R$ containing the set $S$.

\paragraph{}
Let $(R,+,\circ)$ be a ring and $I \neq R$ be an ideal in $R$.

\begin{itemize}[leftmargin=15pt]
\item[1.]
$I$ is called a \textit{maximal ideal}, 
if it is not contained in a larger ideal besides $R$ itself.
\item[2.]
$I$ is called a \textit{primary ideal}, if the following is true:
If $x,y \xeof R \setminus I$, then $x \circ y \xeof R \setminus I$,
that is,
the set $R \setminus I$ is closed {\wrt} multiplication.
\item[3.]
The set $\operatorname{Spec}(R)$ of all primary ideals in $R$
is called the \textit{spectrum} of $R$.
The set of all maximal ideals in $R$ is called the \textit{maximal spectrum}
of $R$.
\item[4.]
$I$ is called a \textit{principal ideal},
if it is multiplicatively generated by one element,
that is,
there is a $y \xeof R$ with
$I \eq \set{x \xeof R}{x \eq r \cdot y, \, r \xeof R}$ rsp. $I \eq [y]$.
\end{itemize}

\paragraph{} %- principal ideal rings
Let $R$ a commutative ring without zero divisors,but with unit element.
Then $R$ is called a \textit{principal ideal ring},
if each ideal in $R$ is a principal one.

\begin{itemize}[leftmargin=15pt]
\item[1.]
Each ring $R$ and the mutually existing unit element $\{1\}$ is a ideal.

\item[2.]
Let $I$ be an interval
and $(C(I),+,\cdot)$ be the ring of continuous functions $f \:I \xto \R$.
For $x \xeof I$ let $I_x := \{f \xeof C(I) \: f(x) \eq 0 \}$.
Then $I_x$ is a maximal ideal in $C(I)$.
Moreover,
each maximal ideal in $C(I)$ is a set $I_x$ for an element $x \xeof I$.

\item[3.]
Let $R \eq (\Z,+,\cdot)$ and 
$m \, \Z := \{z \xeof \Z \: z = m \cdot i, \, i \xeof \Z \}$ for $m \xeof \N$.
Then the sets $m \, \Z$ are principal ideals in $\Z$.
The ring $\Z$ is commutative and withour zero divisors.
Since each ideal in $\Z$ is principal, 
$\Z$ proves to be a principal ideal ring.
\end{itemize}

\section{Factor Rings}
%------------------------------------------------------------------------------%
Let $(R,+,\circ)$ be a ring and $I \subset R$ be a two-sided ideal. 
Since $I$ forms an Abelian subgroup of $R$,
we can consider the (Abelian) factor group $(R/I,+)$.

\subparagraph{}
For each pair of elements $[x],[y] \xeof R/I$ the multiplicative operation
$[x] \circ [y] := [x \circ y]$ is well-defined.
The ideal $I$ is two-sided in $R$,
hence the multiplication is indepenedent of the representatives
$x \xeof [x]$ and $y \xeof [y]$.
Now,
the triple $(R/I,+,\circ)$ has also the structure of a ring,
so the following definition is justified:
$R/I$ is called the \textit{factor ring} rsp.
the \textit{residual class ring} of $R$ by $I$.

\subparagraph{}
If $R$ is a commutative ring,
then each factor ring $R/I$ is commutative, too,
and if $R$ owns a unit element, 
then for each factor ring $R/I$ a unit element can be constructed, too. 
For $R \eq \Z$ and $m \xeof \N$ the ring $\Z_m := \Z / (m \, \Z)$
is the residual class ring of $\Z$ factorized by the principal ideal $m \, \Z$.
The rings $\Z_m$ constitute commutative rings with unit elements.

\paragraph{}
Like for other algebraic structures,
the notion of \textit{homomorphism} is also defined for rings.
Let $(R,+,\circ)$ and $(S,+,\bullet)$ be rings.
Then a \textit{ring homomorphism} is a group homomorphism $\phi \:R \xto S$
that is compatible with the multiplicative operations in $R$ and $S$
in the sense of
\[
\phi(x \circ y) = \phi(x) \bullet \phi(y), \qquad x,y \in R.
\]

\subparagraph{}
Other concepts like epimorphisms,
monomorphisms, isomorphisms and endomorphisms are defined for rings as well.
The rings form together with their homomorphisms a category.

\subparagraph{}
Let $(R,+,\circ)$ a ring and $I$ a two-sided ideal in $R$.
Then the residual class assignment $\pi\:R \xto R/I$, $x \mapsto [x]$
is a ring homomorphism,
denoted the \textit{canonical epimorphism} of $R$ to $R/I$.

\subparagraph{}
The \textit{kernel} $\kernel{\phi}$ of a ring homomorphism $\phi\:R_1 \xto R_2$
is the set of all $x \xeof R_1$ with $\phi(x) \eq 0$.
If $\phi$ is surjective,
then $\kernel{\phi}$ forms a two-sided ideal in $R_1$.
In this case the ring $R_1$ can be factorized by $\kernel{\phi}$,
and the following theorem is valid:

\begin{theorem}
Let $R_1$, $R_2$ be rings and $\phi\:R_1 \xto R_2$ be a surjective
ring homomorphism.
Then the ideals $R_1 / \kernel{\phi}$ and $R_2$ are isomorphic to each other.
\end{theorem}

\paragraph{}
In some cases a variant of the concept of ring homomorphisms is meaningful,
which is denoted \textit{derivations} over a given ring:
Let $(R_1,+,\circ)$ and $(R_2,+,\bullet)$ be rings.
Then a \textit{derivative} is a group homomorphism $\phi\:R_1 \xto R_2$,
if it is compatible with the ring multiplications,
that is
\[
\phi(x \circ y) \= \phi(x) \bullet y \, + \, x \bullet \phi(y),
\quad \fa x,y \in R_1.
\]

\subparagraph{}
Look at the following example:
let $R := C^1(\R)$ be the ring of countinuously differentiable functions
$f\:\R \xto \R$
and $S := C(\R)$ be the ring of continuous functions $f\:\R \xto \R$.
Then the differential operator $\phi \: R \xto S$ 
defined by $\phi(f) := f'$ for $f \xeof C^1(\R)$ forms a derivation
from $R$ to $S$,
since the product rule $(fg)' \eq f'g \, + \, fg'$ holds.

\section{Polynomial Rings}
%------------------------------------------------------------------------------%
For a commutative ring $R$ with unit element
let $R_0$ be the set of \textit{terminating sequences} $f\:\N \xto R$.
It should be clear
that for two sequences $p,q \xeof R_0$ the addition
$p + q \xeof R_0$ is defined component-wise.
Moreover,
theres is a multiplicative operation $\bullet$ in $R_0$,
which is close to that of convolution:
\[
(p \bullet q)_k \eqdef (\sum_{i+j=k} p_i \cdot q_j), \quad k \in \N.
\]
Hence a ring structure on $R_0$ is defined by these operations.

\subparagraph{}
The triple $(R_0,+,\bullet)$ is called the \textit{polynomial ring}
associated with $R$,
and each element $p \xeof R_0$ is called a \textit{polynomial}.
The unit elements of a polynomial ring are given by the sequences
\[
\0 : = (0,0,\ldots), \, \mbox{ and } \, 1 := (1,0,\ldots).
\]
Morover,
if a poynomial ring owns a generating element $x := (0,1,0,\ldots)$,
{\ie} $R_0 \eq [x]$, 
then each other element $p$ of the polynomial ring can be represented by
\[
p = a_0 + a_1 \circ x + a_2 + x^2 + \ldots + a_n \circ x^n,
\]
which coincides with the usual notations for polynomials.
We call the elements $a_i$ the \textit{coefficients} of this polynomial.

\subparagraph{}
By the \textit{degree} of a polynomial $p$ is meant the number
$\grad{p} := \max \set{a_k}{k \xeof \N_0}$.
The following computation rules are valid for all $p,q \xeof R_0$:
\[
\begin{cases}[1.4]
\quad \grad{p+q}         \, \leq \, \max \{ \grad{p}, \grad{q} \},\\
\quad \grad{p \bullet q} \, \leq \, \grad{p} + \grad{q}.\\
\end{cases}
\]
It is of interest to know
that by inserting a ring element $x \xeof R$ into a polynomal $p \xeof R_0$
we obtain a ring homomorphism
\[
\phi_x\:R_0 \to R, \quad p \mapsto p(x),
\]
which is often called the \textit{insertion homomorphism}.

\section{Fields}  
%------------------------------------------------------------------------------%
By a \textit{field} is meant a ring $(R,+,\circ)$,
for which all its elements besides $\0 \xeof R$ own an inverse element.
Let $(K,+,\circ)$ a ring with unit element $\1$.

\begin{itemize}[leftmargin=15pt]
\item[1.]
$K$ is called a \textit{skew field},
if for each element $x \xeof K \setminus \{0\}$ there is $x^{-1} \xeof K$ 
such that $x \circ x^{-1} \eq x^{-1} \circ x \eq \1$.
\item[2.]
In addition,
if $K$ is commutative, then $K$ is denoted a \textit{field}.
\end{itemize}

\subparagraph{}
If $K$ is a skew field,
then the pair $(K \setminus \{0\},\circ)$ forms a group.
The sets $\Q, \R, \C$ of rational, real and complex number
constitute also fields,
but the set of \textit{quaternions} however is only a skew field.

\paragraph{}
We next answer the question,
whether for a given ring $R$ and ideal $I$ in $R$
the factor ring $R/I$ forms a field.
The following result is valid:
let $(R,+,\circ)$ be a commutative ring with unit element 
and $I$ an ideal in $R$. 
Then $R/I$ is a field \iff
$I$ is a \textit{maximal} ideal.

\subparagraph{}
Let $K$ be a field and $\1 \xeof K$ be the unit element {\wrt} multiplication.
Then the smallst integer $\chi(K)$ having the property
\[
\underbrace{\1 \,+\, \ldots \,+\, \1}_{\chi(K)} \= 0
\]
is called the \textit{characteristic} of $K$.
If no such number exists,
then the field $K$ is said to have characteristic $\chi(K) \eq 0$.
An example of a field with $\chi(K) \eq 0$ is the field $\Q$ of rational numbers.

\subparagraph{}
Given an integrity ring $R$,
one can construct a field,
called the \textit{quotient field} of $R$:
First define an equivalence relation on the cartesian product
$\setM := R \times R \setminus \{0\}$ according to
\[
(a,b) \sim (c,d) \, \Leftrightarrow : \, a \circ d = b \circ c.
\]
We denote the equivalence class of $(a,b)$ briefly by $\frac{a}{b}$.
The equivalence classes of $\sim$ form the basic set $Q$
of the quotient field:
\[
Q \eqdef \setM / \sim \= \{ \frac{a}{b} \: (a,b) \xeof \setM \}.
\]
But addition and multiplication on $Q$ is still required:
\[
\frac{a}{b} + \frac{c}{d} := \frac{(a \circ d) + (c \circ b)}{b \circ d},
\quad \quad
\frac{a}{b} \bullet \frac{c}{d} := \frac{a \circ c}{b \circ d}.
\]
The reduction rule in integrity rings implies
that these operations are independent of the representants
of the equivalence classes.

\subparagraph{}
There is still the task to identify the unit elements {\wrt}
addition and multiplication:
\[
0 \eqdef \frac{0}{1},
\quad
1 \eqdef \frac{1}{1}.
\]
Then the triple $(Q,+,\bullet)$ forms a field,
in which the original ring $R$ can be embedded by an injective mapping $\iota$,
that is,
there is a monomorphism $\iota \:R \xto Q$ given by $x \mapsto \frac{x}{1}$.

\subparagraph{}
The most well-known example for this construction is the definition
of the field $\Q$ of rational numbers --
it is the quotient field of the integrity ring $\Z$.

%------------------------------------------------------------------------------%
\fancyhead[EC] {\small{II. \textsc{Linear Spaces and Mappings}}}
\fancyhead[OC] {\small{II. \textsc{Linear Spaces and Mappings}}}

\part{Linear Spaces and Mappings}
%------------------------------------------------------------------------------%
An Abelian group $(\vecV,+)$ is called a \textit{$\F$-linear space},
if there is an outer multiplication $\cdot \: \F \times \vecV \xto \vecV$,
which fulfills for all $\lambda,\mu \xeof \F$ and $x,y \xeof \vecV$
\[
\begin{cases}[1.2]
\quad \lambda \cdot (x + y)       = \lambda \cdot x + \lambda \cdot y,\\
\quad (\lambda + \mu) \cdot x     = \lambda \cdot x + \mu \cdot x,\\
\quad \lambda \cdot (\mu \cdot x) = (\lambda \mu) \cdot x,\\
\quad 1 \cdot x                   = x.
\end{cases}
\]
The mapping $(\lambda,x) \mapsto \lambda \cdot x$
(briefly denoted $\lambda x$)
is called \textit{scalar multiplication}.
If $\F$ is the field of real or complex numbers,
then $\vecV$ is called a \textit{real} or \textit{complex} linear space,
respectively.
The previous conditions imply $0 \cdot x \eq \0$ for each $x \xeof \vecV$.

\section{Linear Subspaces}
%------------------------------------------------------------------------------%
A subset $\vecU \subseteq \vecV$ is called a \textit{subspace} of $\vecV$,
if $\vecU$ is a subgroup of $\vecV$ and
$\lambda \cdot x \xeof \vecU$ holds for all $x \xeof \vecU, \lambda \xeof \F$,
so $\vecU$ forms a linear space in its own right.

\subparagraph{}
One can verify
that the intersection of arbitrarily many subspaces of a given linear space
is a subspace again.
So for each non-empty set $\genG \subset \vecV$
it is possible to define (like in case of groups and rings)
a smallest subspace $[\genG]$ containing $\genG$,
which is said to be generated by $\genG$:
the set $[\genG]$ is the intersection of all subspaces $U \supseteq \genG$
and forms the smallest subspace in $\vecV$ containing the set $\genG$.
The subspace $[\genG]$ is called the \textit{span} of the set $\genG$.
Any vector $x \xeof \vecV$ is an element of $[\genG]$, if and only if
$x$ can be represented as \textit{linear combination} of vectors in $\genG$,
that is, there are vectors $x_1,\ldots,x_n \xeof \genG$ and
coefficients $\xi_1,\ldots,\xi_n \xeof \F$ such that
$x \eq \xi_1 x_1 + \xi_2 x_2 + \ldots + \xi_n x_n$.

\paragraph{} %- Hamel bases
One of the main concepts in the theory of linear spaces
is that of a \textit{Hamel base}.
Preliminary,
a non-empty set $B \subset \vecV$ is called \textit{linearly independent},
if for each selection $\{e_1,\ldots,e_n\}$ of vectors in $B$
the condition
$\xi_1 \, e_1 + \xi_2 \, e_2 + \ldots + \xi_n \, e_n \eq \0$ implies
$\xi_i \eq 0$ for $i \eq 1,\ldots,n$.
Note
that the null vector cannot be an element of any linearly independent set.

\subparagraph{}
A \textit{Hamel base} of a linear space $\vecV$ is
a linearly independent subset $\basH$,
which generates the entire linear space,
{\ie} $[\basH] \eq \vecV$.
In this case the set $\basH$ cannot be extended to a larger linearly
independent set by adding further vectors $\0 \neq v \notin \basH$.

\paragraph{}
The question arises,
whether each linear space has a Hamel basis.
Indeed,
this is true if the \textit{axiom of choice} in set theory is provided:
{\em
assume the axiom of choice to be valid,
then each linear space $\vecV \neq \{\0\}$ has a Hamel basis.}
Induced by set inclusion,
there is a natural partial order on the family
of linearly independent subsets in $\vecV$,
since each chain of linearly independent sets is bounded
by the linear space $\vecV$ itself.
Due to \textit{Zorn's lemma},
which is equivalent to the axiom of choice,
there is a maximal linearly independent set $\basH$
being a Hamel base in $\vecV$.

\paragraph{}
Two Hamel bases $\basH,\basH'$ of a linear space $\vecV$
have always identical cardinalities,
that is, there is a bijective mapping $\phi \:\basH \xto \basH'$.
By this reason,
the linear space $\vecV$ is called \textit{free of rank $\abs{\basB} $},
and we can speak of the \textit{dimension} of a linear space $\vecV$,
which is defined to be the cardinality of any Hamel base in $\vecV$
(we will assign the dimension $0$ to the linear space
consisting of the null vector only).

\subparagraph{}
If a Hamel base of $\vecV$ consists of a finite number of elements,
then the linear space $\vecV$ is called \textit{finite-dimensional}
respectively \textit{n-dimensional},
otherwise the space is said to be \textit{infinite-dimensional}.

\section{Convexity}
%------------------------------------------------------------------------------%
A subset $\convC$ of a linear space $\vecV$ is called \textit{convex},
if for any selection $v_1, \ldots, v_n$ of elements in $\convC$
and coefficients $\lambda_1,\ldots,\lambda_n \xeof \F$ satisfying
$0 \leq \lambda_i \leq 1$ and
$\sum_{i=1}^n \lambda_i \eq 1$
the \textit{convex combination}
\[
C[v_1,\ldots,v_n;\lambda_1,\ldots,\lambda_n]
\eqdef
\lambda_1 \, v_1 + \ldots + \lambda_n \, v_n
\]
is again an element of $\convC$.
The \textit{line segment} $[x,y]$ for $x,y \xeof \vecV$ is
a set of vectors, for which each element $v \xeof [x,y]$
can be represented by the equation $v \eq \lambda x + (1-\lambda) y$
for some $\lambda \xeof [0,1]$.
In addition, the {\em inner} line segment $(x,y)$ consists of those vectors,
which satisfy this equation for $\lambda \xeof (0,1)$.
So an alternative definition for the concept of convexity is possible:
a subset $\convC$ is convex,
if and only if
$(x,y) \subset \convC$ for each pair of vectors $x,y \xeof \convC$.

\subparagraph{}
Notice
that for a linear space $\vecV$ the empty set and the space itself
are always convex.
For example,
the convex subsets of the real line are exactly the intervals
of finite or infinite type.

\paragraph{}
Convex sets have the property,
that intersections of arbitrarily many convex sets are convex again.
This property enables us to declare the \textit{convex hull} $\convC(\setM)$
for each subset $\setM$ of $\vecV$,
which is defined as the intersection of all convex sets
$\setM' \supseteq \setM$.
The operation, which maps each subset of a linear space to its closed hull,
is a hull operator.
Thus we have $\setM \subseteq \convC(\setM)$
and $\convC(\convC(\setM)) \eq \convC(\setM)$.

\section{Linear Mappings}
%------------------------------------------------------------------------------%
Let $\vecV$ and $\vecW$ be arbitrary $\F$-linear spaces.
Then each mapping $\linT \:\vecV \xto \vecW$ satisfying the condition
\[
\linT(\lambda \, x + \mu \, y) \= \lambda \, \linT(x) + \mu \, \linT(y),
\quad (x,y \in \vecV, \; \lambda, \mu \in \F)
\]
is called \textit{linear trans\-formation} or
a \textit{homo\-morphism} from $\vecV$ to $\vecW$;
often it is called a \textit{linear operator}.
We denote $\LO{\vecV,\vecW}$ the set of all
linear trans\-formations between $\vecV$ and $\vecW$.
Obviously,
with the pointwise defined operations
$\linT_1 + \linT_2$ and $\lambda \, \linT$ ($\lambda \xeof \F$),
the set $\LO{\vecV,\vecW}$ forms a $\F$-linear space again.

\paragraph{Kernels, Ranges and Cokernels}
To each linear trans\-formation $\linT \:\vecV \xto \vecW$
there are related three further linear spaces,
namely the \textit{kernel}, the \textit{range}
and the \textit{cokernel} of $\linT$.
The \textit{kernel} $\kernel{\linT}$ of
$\linT$ is the set of all vectors $x \xeof \vecV$ satisfying $\linT x \eq \0$,
which turns out to be a subspace of $\vecV$.
Similarly,
the \textit{range} $\range{\linT} := \set{\linT x \xeof \vecW}{x \xeof \vecV}$
is a subspace of $\vecW$.

\subparagraph{}
To define the cokernel of $\linT$,
we consider the factor group $V/T := \vecV / \kernel{\linT}$.
This group is well-defined due to the fact
that the subgroup $\kernel{\linT}$ is \textit{normal} in $\vecV$.
We still have to endow the group $V/T$ with a linear structure:
let
$\circ \: \F \times V/T \xto V/T$
be the outer multiplication on $V/T$ defined by
\[
\lambda \circ [x] \eqdef
[\lambda x] \quad \mbox{for } [x] \in V/T.
\]
This operation is independent of the representants in $[x]$,
thus $V/T$ is even a $\F$-linear space,
called the \textit{cokernel} of $\vecV$ {\wrt} $\linT$.

\paragraph{}
A vector $x \neq \0$ is called an \textit{eigenvector} to the
\textit{eigenvalue} $\lambda \xeof \F$ {\wrt} the linear transformation $\linT$,
if the equation $\linT x \eq \lambda \, x$ holds.
It can be easily shown
that the \textit{eigenspace} $E_{\lambda}(\linT)$ 
consisting of all eigenvectors to an eigenvalue $\lambda$ 
forms a subspace of $\vecV$, too.

\section{Linear Functionals}
%------------------------------------------------------------------------------%
As a special case of a linear trans\-formation,
a \textit{linear functional} is a mapping $f \:\vecV \xto \F$
satisfying the condition
\[
f(\lambda x + \mu y) \= \lambda f(x) + \mu f(y),
\quad \,
(x,y \in \vecV, \; \lambda,\mu \in \F).
\]
The set $\vecV^* := \LO{\vecV,\F}$ consisting of all linear functionals
on $\vecV$ is clearly a linear space again,
called the \textit{algebraic dual space} of $\vecV$.
Notice
that if $\vecU_1,\vecU_2$ are subspaces of $\vecV$ with
$\vecU_1 \subset \vecU_2$,
then each linear functional $f \xeof \vecU_2^*$ induces
a linear functional $f' \xeof \vecU_1^*$ by restriction to $\vecU_1$,
that is, we have $\vecU_2^* \subset \vecU_1^*$.
This means
that it is possible to increase the dual space by decreasing the base space.

\paragraph{}
Furthermore,
the \textit{bidual space} $\vecV^{**}$ of $\vecV$ is the linear space
consisting of all linear functionals $\Phi \:\vecV^* \xto \F$.
In particular,
each vector $x \xeof \vecV$ can be sent to an element $x^{**}$
of the bidual space $\vecV^{**}$ by the mapping
$\iota \: x \mapsto x^{**}(\Phi) := \Phi(x)$ for $\Phi \xeof \vecV^*$.
This mapping is linear,
hence its image forms a subspace of $\vecV^{**}$.
It is remarkable
that the mapping $\iota$ may be defined without any specification
of a Hamel base in $\vecV$.
Thus $\iota$ is said to be a \textit{natural} linear space homomorphism
from $\vecV$ into its bidual space $\vecV^{**}$.

\subparagraph{}
The embedding $\iota \:\vecV \xto \vecV^{**}$ is always injective.
If the linear space $\vecV$ is finite-dimensional,
then $\iota$ is actually surjective,
and in this case
the spaces $\vecV$ and $\vecV^{**}$ are canonically isomorphic to each other.
However,
in case of an infinite-dimensional linear space
no such assertion can be done without further assumptions.

\paragraph{}
Let $\vecV, \vecW$ be vecor spaces and $\vecV^*, \vecW^*$ their
dual spaces.
Given a linear trans\-formation $\linT \:\vecV \xto \vecW$,
each linear functional $f \xeof \vecW^*$ gives rise to
a linear functional $g \xeof \vecV^*$ by the composition $g:= f \circ \linT\,$.
The assignment $f \mapsto g$ induces a mapping
$\linT^* \:\vecW^* \xto \vecV^*$,
called the \textit{adjoint trans\-formation} of $\linT$.
The adjoint is also linear.
Moreover,
if $\linT$ is bijective, then $\linT^*$ is so.

\begin{center}
\scalebox{0.9}{
\begin{tikzpicture}[node distance=5.0cm]

\node     (v) at (0,0)  {$V$};
\node [right of=v]  (w)     {$W$};

\draw [-stealth] (v)  -- node [midway, above]    {$\linT$}  (w);
\path [stealth-] (v) -- node [midway, below, yshift=-2.5cm] (k) {$\F$} (w);

\draw [-stealth] (v)  -- node [midway, right, yshift=0.2cm] (g') {$g$} (k);
\draw [-stealth] (w)  -- node [midway, left, yshift=0.2cm]  (f)  {$f$} (k);
\draw [-stealth] (f)  -- node [midway, above] {$\linT^*$}    (g');

\end{tikzpicture}
}

\footnotesize{Figure 1: definition of the adjoint transformation.}
\end{center}

\section{Sesquilinear Forms}
%------------------------------------------------------------------------------%
Let $\vecV$ be a complex linear space.
We denote a \textit{sesquilinear form} on $\vecV$
as a mapping $\sql \:\vecV \times \vecV \xto \C$,
which is linear in the first and conjugate linear in the second argument,
that is,
for all $x,y,z \xeof \vecV$ and $\lambda,\mu \xeof \C$ the following rules hold:
\[
\begin{cases}[1.2]
\quad \sql(x+y,z) \= \sql(x,z) \, + \, \sql(y,z)\,,\\
\quad \sql(x,y+z) \= \sql(x,y) \, + \, \sql(x,z)\,,\\
\quad \sql(\lambda x,\mu y) \=
      \lambda \overline{\mu} \, \sql(x,y)\,.\\
\end{cases}
\]
If the underlaying linear space is real,
then $\sql$ is also linear in the second argument and 
thus called a \textit{bilinear form}.
The mapping $x \mapsto \sql[x] := \sql(x,x)$ is called
the \textit{quadratic form} related to $\sql$.

\subparagraph{}
Furthermore,
an \textit{conjugate linear functional} on $\vecV$
is a mapping $f \:\topX \xto \C$ having the properties
\[
f(x+y) = f(x) + f(y)\,, \quad
f(\lambda x) = \overline{\lambda}f(x)
\quad \quad
(x,y \in \topX, \; \lambda \in \C).
\]
Then each sesquilinear form $\sql$ on $\vecV$ can be considered
as a linear map $\linA \:\vecV \xto \vecV^*$,
where $\vecV^*$ is the algebraic dual space of $\vecV$,
that is,
the linear space of all conjugate linear functionals on $\vecV$,
where addition and scalar multiplication are defined pointwise.
Indeed,
given a sesquilinear form $\sql$ on $\vecV$,
each vector $x \xeof \vecV$ gives rise to the conjugate linear functional
\[
f_x(y) := \sql(x,y) \quad (y \in \vecV).
\]
The mapping $\linA \:V \xto V^*$ defined by $\linA \, x := f_x$ is linear
and sometimes called the \textit{representation operator} for $\sql$.
Reversely,
if $\linA \:\vecV \xto \vecV^*$ is linear,
then there is a sesquilinear form $\sql \:\vecV \times \vecV \xto \F$
associated with $\linA$ according to
$
\sql(x,y) := (\linA \,x) y$ for all $x,y \xeof \vecV.
$
Hence sesquilinear forms on a complex linear space $\vecV$
are in one-to-one relationship to linear trans\-formations
between the spaces $\vecV$ and $\vecV^*$.

\paragraph{}
By an \textit{inner product} in a complex linear space $\vecV$
we mean a specific sesquilinear form
$\scalar{\cdot}{\cdot} \:\cross{\vecV}{\vecV} \xto \C$
satisfying the additional conditions for all $x,y \xeof \vecV$:

\begin{itemize}[leftmargin=35pt]
\item[(I-1)]
$\scalar{x}{y} \eq \overline{\scalar{y}{x}}$ \; \textit{(hermitean symmetry)}.

\item[(I-2)]
$\scalar{x}{x} \geq 0$ \; \textit{(positivity)}.

\item[(I-3)]
$\scalar{x}{x} \eq 0$ implies $x \eq 0$ \; \textit{(regularity)}.
\end{itemize}

\subparagraph{}
Linear spaces endowed with an inner product
$\scalar{\cdot}{\cdot}$ are said to be \textit{unitary spaces} or
\textit{inner product spaces}.

\paragraph{}
Due to the existence of the inner product in $\vecV$,
the concept of \textit{orthogonality} in such spaces is really new
in comparison to ordinary linear spaces.
Two vectors $x,y \xeof \vecV$ are called \textit{orthogonal} to each other,
if $\scalar{x}{y} \eq 0$ (notation: $x \bot y$).
A subset $\orthO \subset \vecV$ is called \textit{orthogonal}
or an \textit{orthogonal system},
if any pairs of vectors $x,y \xeof \orthO$ are orthogonal to each other.
It is easy to verify
that an orthogonal system forms a linearly independent set.
For a subset $\setA \subset \vecV$ we introduce the \textit{orthogonal space}
$ \setA^{\bot}$ of $\setA$ by means of
\[
\setA^{\bot} \eqdef
\{ x \in \vecV: \scalar{x}{y} \eq 0 \fa y \in \setA \}.
\]

%------------------------------------------------------------------------------%
\fancyhead[EC] {\small{III. \textsc{Algebras}}}
\fancyhead[OC] {\small{III. \textsc{Algebras}}}

\part{Algebras}
%------------------------------------------------------------------------------%
An \textit{algebra} is a linear space $\algA$,
whose underlying Abelian group has a ring structure
being compatible with the scalar multiplication in a certain way.
More precisely,
let $(\algA,+,\circ)$ be a ring and $\F$ be a field, then
$\algA$ is called a \textit{$\F$-algebra},
if there is an operation $\cdot \:\cross{\F}{\algA} \xto \algA$
such that $(\algA,+)$ is a $\F$-linear space
satisfying the rule
\[
(\lambda \cdot x) \circ y
\=
x \circ (\lambda \cdot y) = \lambda \cdot (x \circ y)
\quad \quad
(\lambda \in \F, \, x,y \in \algA).
\]
An algebra $\algA$ is called \textit{unital},
if it has an identity element $\1 \xeof \algA$ {\wrt} the ring operation $\circ$.
Finally, $\algA$ is called a \textit{commutative} algebra,
if $\algA$ (considered as a ring) is commutative.
For abbreviation one often writes $xy$ instead of $x \circ y$.

\section{Algebra Homomorphisms}
%------------------------------------------------------------------------------%
For some $\F$-algebras $\algA$ and $\algB$,
an \textit{algebra homomorphism} $\Phi \:\algA \xto \algB$
is a linear trans\-formation,
which is compliant with the ring operation, that is
\[
\Phi(x \circ y) \= \Phi(x) \circ \Phi(y) \quad \quad (x,y \in \algA).
\]
A bijective character $\Phi$ between algebras is called
an \textit{algebra isomorphism},
and in this case it can easily be proved
that the inverse mapping $\Phi^{-1}$ is an algebra homomorphism, too.
In the special case,
where $\algB$ is the field $\C$ of complex numbers,
an algebra homomorphism is usually called a \textit{character}.
The existence of an algebra homomorphism between arbitrary algebras 
$\algA$ and $\algB$ is denoted by $\algA \cong \algB$.

\paragraph{}
There is a natural way to endow a non-unital algebra with an identity element.
If $\algA$ is an $\F$-algebra \textit{without} identity,
then an extended algebra \textit{with} identity can be created by
$\algB := \algA \oplus \F$ and $\1 := (\0,1)$,
where $\algA \oplus \F$ is the cartesian product of $\algA$ and $\F$
endowed with algebra operations componentwise defined as follows:
\[
\begin{cases}[1.2]
\quad (x,\lambda) + (y,\mu) \eqdef (x+y, \lambda + \mu)\\
\quad (x,\lambda) (y,\mu)   \eqdef (xy + \mu x + \lambda y, \lambda \mu) \\
\quad \mu (x,\lambda)       \eqdef (\mu x, \mu \lambda)\\
\end{cases}
\]
for all $x,y \xeof \algA$ and $\lambda,\mu \xeof \F$.
Notice
that he original algebra $\algA$ is a two-sided ideal within
the extended algebra $\algB$,
and the factor algebra $\algB/\algA$ is isomorphic to the field $\F$
(considered as a $\F$-algebra): $\algB/\algA \cong \F$.

\section{Subalgebras}
%------------------------------------------------------------------------------%
Let $(\algA,+,\circ)$ be a $\F$-algebra.
A subset $\algU \subset \algA$ is called a \textit{subalgebra} of $\algA$,
if $\algU$ is a subspace of $\algA$ and a subring of $(\algA,+,\circ)$ together.
Similar to subspaces of linear spaces,
the intersection of arbitrarily many subalgebras of $\algA$ is a subalgebra
of $\algA$, too.
If $\genG \subset \algA$ is a non-empty subset,
then the intersection $[\genG]$ of all subalgebras $\algA'$ with the property
$\algA' \supseteq \genG$ is called the subalgebra \textit{generated by $\genG$}.
Finally,
if for any subset $\genG$ the equality $[\genG] \eq \algA$ holds,
then $\genG$ is called a \textit{generator} of $\algA$.

\section{Group Algebras}
%------------------------------------------------------------------------------%
For every group $\grpG$ it is possible to construct
an algebra in a natural way
by performing the formal linear combinations
\[
\alpha_1 \cdot g_1 \, + \, \alpha_2 \cdot g_2 + \ldots + \alpha_n \cdot g_n
\]
with coefficients $\alpha_i \xeof \R$.
The set $A_{\R}(\grpG)$ of these linear combinations constitutes
a $\R$-linear space.

\subparagraph{}
An additional ring structure on $A_{\R}(\grpG)$ is given according to
\[
\sum_{i=1}^n \alpha_i \cdot g_i \, \bullet \,
\sum_{j=1}^m \beta_j  \cdot g_j
\eqdef
\sum_{i,j=1}^{i=n,j=m} \alpha_i \beta_j \, \cdot \, g_i g_j.
\]
The structure $(A_{\R}(\grpG), + , \bullet)$ is called
the \textit{group algebra} of $\grpG$ over $\R$.
The construction of the group algebra can be done basically for each ring $R$.

\section{Ideals}
%------------------------------------------------------------------------------%
A subset $\idI \subset \algA$ is called an \textit{ideal} in $\algA$,
if $\idI$ is a subspace in the linear space $(\algA,+)$
and an ideal in the ring $(\algA,+,\circ)$ together.
Because the intersection of arbitrarily many ideals in $\algA$
is an ideal in $\algA$ again,
we can assign to each non-empty subset $\genG \subset \algA$
the minimal ideal $[\genG]$, which is called the
\textit{ideal generated by $\genG$}.

\subparagraph{}
Let $\algA / \idI$ be the quotient linear space of an ideal $\idI$ in $\algA$.
Then $\algA/\idI$ is also a subring in $\algA$,
so it is possible to factorize each algebra $\algA$
by an ideal $\idI$ in $\algA$.
The set $\algA/\idI$ is a $\F$-Algebra again
and is called the \textit{factor algebra} of $\algA$ by $\idI$.

%------------------------------------------------------------------------------%
%                               END OF DOCUMENT                                %
%------------------------------------------------------------------------------%
\end{document}
